% ===================================================================
% Appendix: matrix-analytic lemmas used in the convexity proof
% ===================================================================
\section{Matrix-analytic lemmas}
\label{app:matrix}

This appendix collects the two matrix-analysis facts used in the proof of
Proposition~\ref{prop:convex}. Both are classical; we include proofs so that the
convexity argument is self-contained, and because the Schur-complement form in
Lemma~\ref{lem:trinv} is also what yields the semidefinite representation of the
A-term referred to in Section~\ref{sec:computation}.

\begin{lemma}[Convexity of the trace of the inverse]
\label{lem:trinv}
The map $\Phi:\Sym^{K}_{++}\to\R$, $\Phi(X)=\tr(X^{-1})$, is convex. Moreover
$\Phi$ admits the semidefinite representation
\begin{equation}
  \tr(X^{-1})\le t
  \quad\Longleftrightarrow\quad
  \exists\,Z\in\Sym^{K}:\ \ \tr Z\le t,\ \
  \begin{bmatrix} Z & I\\ I & X\end{bmatrix}\succeq0 .
  \label{eq:schur}
\end{equation}
\end{lemma}

\begin{proof}
\emph{Convexity.} It suffices to show that $\Phi$ is convex along every line
segment in $\Sym^{K}_{++}$. Fix $X\succ0$ and $H\in\Sym^{K}$, and let
$g(s)=\Phi(X+sH)=\tr\big((X+sH)^{-1}\big)$ for $s$ in a neighbourhood of $0$
small enough that $X+sH\succ0$. Writing $A_s=(X+sH)^{-1}\succ0$ and
differentiating the identity $(X+sH)A_s=I$ twice gives
\[
  \frac{d}{ds}A_s=-A_sHA_s,
  \qquad
  \frac{d^{2}}{ds^{2}}A_s=2A_sHA_sHA_s ,
\]
so that $g''(s)=2\tr\big(A_sHA_sHA_s\big)$. Put $C=A_s^{1/2}HA_s^{1/2}$, which is
symmetric. Then
\[
  A_sHA_sHA_s
  =A_s^{1/2}\big(A_s^{1/2}HA_s^{1/2}\big)\big(A_s^{1/2}HA_s^{1/2}\big)A_s^{1/2}
  =A_s^{1/2}C^{2}A_s^{1/2},
\]
whence $g''(s)=2\tr\big(A_s^{1/2}C^{2}A_s^{1/2}\big)\ge0$, because
$A_s^{1/2}C^{2}A_s^{1/2}$ is positive semidefinite (it is of the form $R^\top R$
with $R=CA_s^{1/2}$). Thus $g$ is convex for every choice of $X$ and $H$, and
therefore $\Phi$ is convex.

\emph{Semidefinite representation.} For $X\succ0$, the Schur complement
criterion states that
$\left[\begin{smallmatrix} Z & I\\ I & X\end{smallmatrix}\right]\succeq0$
holds if and only if $Z-IX^{-1}I=Z-X^{-1}\succeq0$. Hence the right-hand side of
\eqref{eq:schur} asserts the existence of $Z\succeq X^{-1}$ with $\tr Z\le t$. If
such a $Z$ exists then $\tr(X^{-1})\le\tr Z\le t$, since $Z-X^{-1}\succeq0$
implies $\tr(Z-X^{-1})\ge0$. Conversely if $\tr(X^{-1})\le t$ we may take
$Z=X^{-1}$. The set described on the right is defined by a linear matrix
inequality in $(X,Z)$ together with a linear inequality, hence is convex; its
projection onto $(X,t)$ is the epigraph of $\Phi$, which is therefore convex,
giving a second proof of the first claim.
\end{proof}

\begin{lemma}[Concavity of the log-determinant]
\label{lem:logdet}
The map $\Psi:\Sym^{K}_{++}\to\R$, $\Psi(X)=\log\det X$, is concave, and strictly
concave on $\Sym^{K}_{++}$.
\end{lemma}

\begin{proof}
Fix $X\succ0$ and $H\in\Sym^{K}$ with $H\neq0$, and set
$g(s)=\log\det(X+sH)$ for $s$ near $0$. Factoring,
\[
  \det(X+sH)=\det\!\big(X^{1/2}\big(I+sX^{-1/2}HX^{-1/2}\big)X^{1/2}\big)
  =\det X\cdot\det\!\big(I+sW\big),
  \qquad W:=X^{-1/2}HX^{-1/2},
\]
where $W$ is symmetric with real eigenvalues $\omega_1,\dots,\omega_K$. Hence
\[
  g(s)=\log\det X+\sum_{k=1}^{K}\log\big(1+s\omega_k\big),
  \qquad
  g''(s)=-\sum_{k=1}^{K}\frac{\omega_k^{2}}{(1+s\omega_k)^{2}}\;\le\;0 .
\]
So $g$ is concave for every $X$ and $H$, and $\Psi$ is concave. Moreover
$g''(s)=0$ would force $\omega_k=0$ for all $k$, i.e.\ $W=0$ and hence $H=0$,
contrary to assumption; therefore $g''<0$ and $\Psi$ is strictly concave.
\end{proof}

\begin{remark}
Lemma~\ref{lem:logdet} gives strict concavity of $\log\det$ \emph{as a function
of the matrix}. Strict convexity of $\w\mapsto-\log\det\M$ requires in addition
that the linear map $\w\mapsto\M$ be injective, which is the content of
Theorem~\ref{thm:uniqueness}(iv) and Remark~\ref{rem:injective}.
\end{remark}
